% =============================================================================
% Archived from DM.tex on 2026-05-23.
% Source: original \subsection{Quantitative propagator-mismatch metrics}
%         under sec:diagnostics (formerly DM.tex lines 881-928).
%
% Labels defined here:
%   sec:prop_mismatch_metrics
%   eq:propagator
%   eq:per_mode_kl
%   eq:prop_zscore
%   eq:ps_weighted_delta
%   eq:Dtot
%   eq:L1pw
%   eq:Rvar
%
% Surviving references to these labels remain inside the CelebA / FID
% appendix of DM.tex (search for sec:prop_mismatch_metrics, eq:Dtot, etc.)
% and will report as undefined unless this file is \input{...} back in or
% the appendix is updated.
% =============================================================================

\subsection{Quantitative propagator-mismatch metrics}
\label{sec:prop_mismatch_metrics}

The radially averaged momentum-space propagator
\begin{equation}
G(|p|) = \bigl\langle|\tilde\phi(p)|^2\bigr\rangle_{\text{radial}},
\label{eq:propagator}
\end{equation}
with $\tilde\phi(p)$ the discrete Fourier transform of a configuration, encodes the two-point structure of the ensemble. To quantify how well the diffusion-model propagator $G_{\text{DM}}(|p|)$ matches the reference $G_{\text{train}}(|p|)$ across momentum bins, one could compare them via absolute deviation $\Delta G(|p|)=G_{\text{DM}}-G_{\text{train}}$ or relative deviation $\Delta G/G_{\text{train}}$, but neither scales uniformly across the UV/IR range. The absolute deviation is dominated by a handful of low-$k$ bins where $G$ is large, while the relative deviation grows artificially at high $k$ where $G$ is tiny and the mismatch contributes negligibly to the field variance. We instead use three complementary metrics, each with a sharp probabilistic or statistical interpretation.

\paragraph{(i) Per-mode Gaussian KL.}
Under translation invariance, any zero-mean Gaussian field measure diagonalises in momentum with covariance $G(k)$, and the modes become independent. The Kullback--Leibler divergence between the two Gaussian field distributions $p_{\text{train}}=\mathcal{N}(0,G_{\text{train}})$ and $p_{\text{DM}}=\mathcal{N}(0,G_{\text{DM}})$ therefore reduces to an additive sum over momentum modes:
\begin{equation}
D_{\text{KL}}(p_{\text{train}}\|p_{\text{DM}}) = \sum_k D_k,\qquad
D_k = \tfrac{1}{2}\bigl[r_k - 1 - \log r_k\bigr],\quad r_k=\tfrac{G_{\text{train}}(k)}{G_{\text{DM}}(k)}.
\label{eq:per_mode_kl}
\end{equation}
The function $f(r)=r-1-\log r$ vanishes at $r=1$, is strictly convex, grows linearly as $r\to\infty$ (linear penalty for DM underestimating the variance) and only logarithmically as $r\to 0^+$ (weaker penalty for overestimation). For small relative deviations $r=1+\varepsilon$ one has $D_k\simeq\varepsilon^2/4$, so $D_k$ reduces to the usual quadratic relative-deviation measure in the well-trained regime, but is far tamer than $(r-1)^2$ for large mismatches in the UV tail. Because $D_k$ is dimensionless and additive, contributions from different $|k|$-bins (and hence UV vs IR halves) can be compared and summed on equal footing.

Crucially, \eq{eq:per_mode_kl} is the KL between two \emph{Gaussian surrogates} determined by the two-point functions alone. When the underlying field theory is non-Gaussian (as in $\phi^4$ with interactions), the metric captures the second-moment mismatch but is blind to higher-cumulant structure such as the connected four-point function $\langle\phi^4\rangle_{\text{conn}}$; matching the two-point function is a necessary but not sufficient condition for distributional equivalence. The generalisation $D_{\text{KL}}(p\|q_{\text{Gauss}})=D_{\text{KL}}(p_{\text{Gauss}}\|q_{\text{Gauss}})+J(p)$, with $J(p)\ge 0$ the negentropy of the true distribution, makes the missing piece explicit.

\paragraph{(ii) Statistical $z$-score.}
Bootstrap-resampling the configurations yields uncertainty estimates $\sigma_{\text{DM}}(k)$ and $\sigma_{\text{train}}(k)$ on the two propagators. The standardised residual
\begin{equation}
z_k = \frac{G_{\text{DM}}(k)-G_{\text{train}}(k)}{\sqrt{\sigma_{\text{DM}}^2(k)+\sigma_{\text{train}}^2(k)}}
\label{eq:prop_zscore}
\end{equation}
tests whether the per-bin discrepancy is statistically significant given the finite sample size. Bins with $|z_k|\lesssim 2$ are consistent with the reference at the $2\sigma$ level; persistent deviations with $|z_k|\gg 2$ across many bins signal a real bias rather than sampling noise. The root-mean-square $z_{\text{rms}}=\sqrt{\langle z_k^2\rangle}$ provides a single scalar that saturates at $\mathcal{O}(1)$ when the mismatch is at the sample-noise floor and grows without bound when the bias dominates.

\paragraph{(iii) Phase-space weighted propagator difference.}
Expressed in the continuum limit, the field variance decomposes as $\langle\phi^2\rangle=\int\! d^Dk/(2\pi)^D\,G(k)\propto\int_0^\infty\! dk\,k^{D-1}G(k)$, so that each radial momentum shell contributes its propagator weighted by the phase-space measure $k^{D-1}$. Accordingly, the density
\begin{equation}
w_k = k^{D-1}\bigl[G_{\text{DM}}(k)-G_{\text{train}}(k)\bigr]
\label{eq:ps_weighted_delta}
\end{equation}
measures, up to an overall $(2\pi)^{-D}$ factor, how much each momentum shell contributes to the total mismatch in $\langle\phi^2\rangle$. Integrating $w_k$ over a UV or IR band directly quantifies the corresponding contribution to the two-point mismatch in position space, with physical units. Compared to $D_k$ and $z_k$, $w_k$ has a concrete physical interpretation but is not dimensionless, so the three metrics should be viewed as complementary rather than redundant.

\paragraph{Aggregate scalars.}
Each per-mode quantity above also admits a single-scalar summary. These scalars are convenient training diagnostics, one number per checkpoint.
\begin{align}
D_{\mathrm{tot}} &= \sum_k D_k,
  \label{eq:Dtot}\\
L_1^{\mathrm{pw}} &= \sum_k k^{D-1}\bigl|G_{\mathrm{DM}}(k)-G_{\mathrm{train}}(k)\bigr|,
  \label{eq:L1pw}\\
R_{\mathrm{var}} &= \frac{\sum_k G_{\mathrm{DM}}(k)}{\sum_k G_{\mathrm{train}}(k)}.
  \label{eq:Rvar}
\end{align}
$D_{\mathrm{tot}}$ is the Gaussian-surrogate KL of \Eq{eq:per_mode_kl} summed over bins, and inherits its ratio-based scale invariance. $L_1^{\mathrm{pw}}$ is the $L_1$ norm of the phase-space-weighted density $w_k$ of \Eq{eq:ps_weighted_delta}, and by Parseval's relation controls the absolute $\langle\phi^2\rangle$ mismatch in physical units. $R_{\mathrm{var}}$ is the total-variance ratio, tracked directly rather than via $|R_{\mathrm{var}}-1|$ because its signed deviation records whether the model underestimates or overestimates the bulk variance at each training stage. The three scalars carry complementary information. $D_{\mathrm{tot}}$ is insensitive to a uniform rescaling of $G$ that $R_{\mathrm{var}}$ catches immediately. $L_1^{\mathrm{pw}}$ weights mid-$k$ modes most strongly and is the closest of the three to a perceptual summary.


% =============================================================================
% Also archived 2026-05-23: \section{Perceptual benchmark: FID and CelebA
% cross-check} (formerly DM.tex lines 1323-1342). This appendix applies the
% propagator-mismatch metrics defined above to a CelebA-trained NCSN++ image
% model and cross-checks against FID; it depends on
%   sec:prop_mismatch_metrics, eq:per_mode_kl, eq:prop_zscore,
%   eq:ps_weighted_delta, eq:Dtot, eq:Rvar
% above.
%
% Labels defined here:
%   app:fid_celeba
%   eq:fid
%   fig:diag_celeba64
% =============================================================================

\section{Perceptual benchmark: FID and CelebA cross-check}
\label{app:fid_celeba}

For image data we cross-check the spectral diagnostics of \Sec{sec:prop_mismatch_metrics} against the Fréchet Inception Distance (FID)~\cite{Heusel2017FID}. Passing a sample set through the Inception-v3 network yields $d=2048$-dimensional activations. Fitting a Gaussian $\mathcal{N}(\mu, \Sigma)$ to these activations for both the DM and training samples, FID is the Fréchet distance between the two Gaussians,
\begin{equation}
\mathrm{FID} = \|\mu_{\mathrm{DM}}-\mu_{\mathrm{train}}\|^2
  + \mathrm{tr}\!\left[\Sigma_{\mathrm{DM}}+\Sigma_{\mathrm{train}}
    -2\bigl(\Sigma_{\mathrm{DM}}\Sigma_{\mathrm{train}}\bigr)^{1/2}\right],
\label{eq:fid}
\end{equation}
which, like the spectral $D_{\mathrm{tot}}$, is a Gaussian KL up to a monotone rescaling, but in perceptual Inception-feature space rather than in momentum space. Comparing FID across training with the spectral scalars in \Eqs{eq:Dtot}{eq:Rvar} tests how much perceptual and physical-spectrum quality track each other.

\fig{fig:diag_celeba64} applies the three propagator-mismatch diagnostics to an NCSN++ diffusion model trained on CelebA grayscale $64\times64$ images, using log-spaced checkpoints from epoch~1 to epoch~$10{,}000$. Each checkpoint supplies $N=512$ EM-sampled images. The reference $G_{\text{train}}$ is computed from $50{,}000$ training images, so $\sigma_{\text{train}}\ll\sigma_{\text{DM}}$ and the DM sample size dominates the bootstrap uncertainty. All three metrics display the expected monotonic improvement during training: the per-mode KL drops by nearly two orders of magnitude at all $|k|$, the $z$-score envelope contracts toward the $\pm 2\sigma$ band, and the phase-space weighted difference $w_k$ shrinks and changes sign as the model passes from underestimating to slightly overestimating the low-$k$ power. The UV (large-$|k|$) portion dominates the total KL throughout training, consistent with the difficulty of learning high-frequency image modes.

\begin{figure*}[t]
\centering
\includegraphics[width=\textwidth]{diag_3panel_celeba_64_em.pdf}
\caption{Propagator-mismatch metrics for an NCSN++ diffusion model on CelebA grayscale $64\times 64$ images at eight log-spaced training checkpoints. Left: per-mode Gaussian KL $D_k$ of \eq{eq:per_mode_kl}; the reference value $D_k=0$ corresponds to exact propagator matching. Centre: statistical $z$-score of \eq{eq:prop_zscore}; the dashed lines mark the $\pm 2\sigma$ consistency band. Right: phase-space weighted propagator difference $w_k$ of \eq{eq:ps_weighted_delta} with $D=2$; integrating $w_k$ over a momentum band returns the contribution of that band to the $\langle\phi^2\rangle$ mismatch.}
\label{fig:diag_celeba64}
\end{figure*}
